\newpage
\ifdefined\useChapters
\chapter{Descaradamente}
\else
\chapter{}
\fi

Nossa, já são dez horas da noite. Como o tempo voa.  
Desde que Lúcia foi para Nova Iorque, a casa voltou a ser silenciosa demais. Não fico mais acampando no meio do nada, nem fugindo tanto — mas esse sossego é sempre interrompido pela sensação de que tudo está à beira de explodir.

Um barulho vindo da sala me tira do transe.  
O telefone.

Teleporto até ele — ainda não aprendi a simplesmente caminhar quando posso ser rápido.

— Alô?

— Alô, quem tá falando?

A voz é inconfundível, mas meu cérebro inicialmente se recusa a encaixar.

— Então… eu sei que é estranho eu te ligar, mas por favor, não desliga ainda. Sou eu. A Larissa.

Eu seguro o ar por um instante.

— Larissa? Que Larissa?

— A de cabelo azul que tentou te queimar soltando labaredas pelas mãos. Aposto que não teve outra assim.

Reviro os olhos.

— Caramba… você é muito cara de pau. Como conseguiu meu número?

Ela ri — não de alegria, mas de nervosismo mascarado.

— Isso não importa agora. Eu preciso muito da sua ajuda. Não estaria te ligando se tivesse outra saída.  
Eu preciso que você encontre o Crispim pra mim.

— O quê? Encontrar o Crispim? Como assim?

Sinto meus músculos tensionarem.  
Eu tento acessar a leitura mental dela, mas não consigo — está longe demais. Só percebo que há algo escondido, uma intenção que ela segura com força.

— Você quer minha ajuda, não é? — digo. — Então precisa me explicar direito o que está acontecendo. E é melhor fazermos isso ao vivo.

Ela faz uma pausa longa o suficiente para eu sentir que está se preparando para um golpe.

— Tudo bem — responde finalmente. — Eu sei que, mesmo que você queira, não vai conseguir fazer nada comigo.

A ousadia dela sempre foi parte do charme e do perigo.

— Me encontre na praça central — digo. — Lá é sempre cheio de gente.

— Amanhã… depois do almoço. Duas horas.  
Aí eu te explico melhor.  
Até mais.

A ligação cai.

Fico parado encarando o nada, sentindo a mistura de irritação e propósito se acomodando no peito.  
Larissa pedindo minha ajuda.  
Larissa, que tentou me queimar viva.  
Larissa, que nunca confiou em mim — e que sempre foi impulsiva demais para admitir medo.

Libertar Crispim…  
Pode ser emboscada?  
Com certeza.

Mas eu vou.  
Não porque confio, mas porque **não tenho mais medo do que podem fazer comigo**.  
Eu regressei do futuro com habilidades que ninguém ali entende — e isso me coloca sempre dois passos à frente.

\bigskip

Na praça central, em plena terça-feira, estou sentado no banco, observando o movimento. O mundo parece outro depois que aprendi a visualizar as pinturas mentais das pessoas. Cada rosto é uma confusão de cores, angústias, esperanças, batalhas internas.  
Ver isso me faz suportar problemas que antes eu detestava.  
Todo mundo, sem exceção, está travando guerras invisíveis.

Sinto Larissa se aproximando a dois quarteirões de distância.  
Seus pensamentos vêm como fagulhas rápidas, impulsivas.  
E sinto também algo mais.

Ítalo.  
Pedro.

Claro que ela não veio sozinha.

Quando finalmente ela surge no meio da multidão, dá para ver no jeito como ela caminha que não dormiu direito.  
A determinação dela vem sempre com um toque de raiva — raiva do mundo, de Crispim, de si mesma.

— Como eu te prometi, estou aqui — diz ela. — E sozinha.

Minto com calma:

— Estou vendo.

Mas por dentro, sei onde estão os outros.  
E ela não pode saber que eu vejo mais do que ela imagina.  
Minha vantagem é justamente **ninguém saber que posso ler pensamentos**.

Ela respira fundo, força um sorriso nervoso, e começa:

— Quando cheguei na casa do Crispim… encontrei o corpo dele desacordado. E, do lado, estava o garçom… fingindo não saber de nada.

A voz dela treme quase imperceptivelmente.  
Ela quer parecer forte — sempre quis — mas está quebrada.  
E a imagem mental dela é um mosaico:  
vermelho de culpa, azul de esperança, e um preto silencioso de medo.

— Eu sei que você é a última pessoa a quem eu deveria pedir ajuda — continua — depois de tudo que te fiz. Mas eu não tenho mais a quem recorrer. E você… pelo que eu sei… é o único que pode ir atrás dele e trazê-lo de volta.

— E por que você acha que eu vou fazer isso?

Ela engole seco.

— Porque você é uma pessoa boa.  
E sabe que ele não teve culpa de nada.  
Ou talvez porque você queira provar alguma coisa.  
Eu realmente não sei.  
Só sei que… é o certo.  
E não quero deixar alguém preso num limbo entre a vida e a morte.

A sinceridade dela bate em mim como um golpe.  
Não pelo conteúdo — mas pelo medo que vem junto.  
Ela está desesperada e tenta esconder isso por trás de bravatas.

— Quem me garante — digo — que assim que eu libertar o Crispim, vocês não vão voltar a perseguir alguém como fizeram comigo?

Ela balança a cabeça, cansada.

— Eu não quero mais saber de perseguir ninguém.  
Eu só quero…  
Eu só quero que isso acabe.

Mentira.  
Mas uma mentira útil.

Decido jogar o mesmo jogo.

— Vou te ajudar — digo, finalmente. — Vou atrás dele.  
Enquanto isso… vou deixar meu corpo aqui.

Ela arregala os olhos, assustada.

— Só… não faça nada com ele.

Antes que ela responda, desdobro.

Meu corpo cai no banco como uma carcaça vazia.

\bigskip

A escuridão me engole.

Que lugar é esse?  
É um vazio profundo, silencioso, onde não existe parede, chão ou teto — só uma neblina escura que parece se mover como se respirasse.

A única coisa visível é uma luz, bem ao fundo.  
Debaixo dela, alguém sentado.

Crispim.

Ou… o que resta dele.

Caminho em direção à luz, cada passo ecoando no nada.

E, pela primeira vez desde que acordei do coma, sinto medo.

Não do lugar.  
Mas do que vou encontrar.